\chapter{Capstone projects}
\label{ch:capstone}

\begin{quote}
\emph{``The best way to learn generative AI is to build something real. These five projects take you from concept to working application.''}
\end{quote}

This chapter presents five capstone projects. Each project integrates concepts from multiple chapters and produces a complete, deployable application. Choose one (or more) based on your interests.

% ============================================================
\section{Project 1: RAG application on a PDF corpus}

\textbf{Goal:} Build a question-answering system over a collection of PDF documents.

\textbf{Chapters used:} 1 (embeddings), 4 (prompting), 6 (RAG).

\subsection{Specification}

\begin{enumerate}
  \item Load 5--10 PDF documents (research papers, textbooks, reports).
  \item Chunk documents with overlap (500--800 characters).
  \item Store embeddings in ChromaDB with metadata (filename, page number).
  \item Build a LangChain RAG chain with Gemini as the LLM.
  \item Add source citations: each answer must reference the document and page.
  \item Evaluate with 10 test questions where you know the correct answer.
\end{enumerate}

\begin{minted}{python}
# Starter code: RAG with source tracking
from langchain_community.document_loaders import PyPDFDirectoryLoader
from langchain_text_splitters import RecursiveCharacterTextSplitter
from langchain_community.vectorstores import Chroma
from langchain_community.embeddings import HuggingFaceEmbeddings
from langchain_google_genai import ChatGoogleGenerativeAI
from langchain_core.prompts import ChatPromptTemplate
from langchain_core.runnables import RunnablePassthrough
from langchain_core.output_parsers import StrOutputParser

# Load all PDFs from a directory
loader = PyPDFDirectoryLoader("./pdf_corpus/")
docs = loader.load()
print(f"Loaded {len(docs)} pages from PDFs")

# Split
splitter = RecursiveCharacterTextSplitter(chunk_size=600, chunk_overlap=80)
chunks = splitter.split_documents(docs)

# Embed and store
embeddings = HuggingFaceEmbeddings(model_name="all-MiniLM-L6-v2")
vectorstore = Chroma.from_documents(chunks, embeddings,
                                     persist_directory="./capstone_rag_db")
retriever = vectorstore.as_retriever(search_kwargs={"k": 5})

# RAG chain with sources
llm = ChatGoogleGenerativeAI(model="gemini-2.0-flash", temperature=0.2)

template = ChatPromptTemplate.from_template(
    """Answer the question based on the context below. Include source
references (document name, page number) for each claim.

Context:
{context}

Question: {question}

Answer (with sources):"""
)

def format_docs_with_sources(docs):
    formatted = []
    for d in docs:
        source = d.metadata.get("source", "unknown")
        page = d.metadata.get("page", "?")
        formatted.append(f"[{source}, p.{page}]\n{d.page_content}")
    return "\n\n".join(formatted)

chain = (
    {"context": retriever | format_docs_with_sources,
     "question": RunnablePassthrough()}
    | template | llm | StrOutputParser()
)

answer = chain.invoke("What are the main findings of the study?")
print(answer)
\end{minted}

\subsection{Deliverables}

\begin{itemize}
  \item Working Jupyter notebook with full pipeline.
  \item Evaluation table: 10 questions, expected answers, generated answers, correctness rating (1--5).
  \item 200-word reflection on RAG limitations encountered.
\end{itemize}

% ============================================================
\section{Project 2: Fine-tune a domain chatbot}

\textbf{Goal:} Fine-tune TinyLlama on a custom instruction dataset to create a domain-specific assistant.

\textbf{Chapters used:} 1 (tokenization), 3 (generation), 5 (fine-tuning).

\subsection{Specification}

\begin{enumerate}
  \item Choose a domain (e.g., cooking, African history, Python programming).
  \item Create 200+ instruction-response pairs (mix of hand-written and curated).
  \item Fine-tune TinyLlama with QLoRA ($r=16$, 3 epochs).
  \item Compare base model vs fine-tuned model on 20 test questions.
  \item Measure perplexity before and after fine-tuning.
\end{enumerate}

\begin{minted}{python}
# Dataset creation helper
import json

def create_instruction_dataset(topic, num_examples=50):
    """Use an LLM to help generate instruction-response pairs."""
    from groq import Groq
    client = Groq()

    pairs = []
    prompt = f"""Generate {num_examples} diverse instruction-response pairs
about {topic}. Format as JSON array with keys: instruction, response.
Each response should be 2-4 sentences. Be accurate and educational."""

    response = client.chat.completions.create(
        model="llama-3.1-8b-instant",
        messages=[{"role": "user", "content": prompt}],
        temperature=0.8, max_tokens=4096,
    )

    try:
        pairs = json.loads(response.choices[0].message.content)
    except json.JSONDecodeError:
        print("Failed to parse JSON. Manual cleanup needed.")

    return pairs

# Generate and save
pairs = create_instruction_dataset("West African cuisine", 50)
with open("cooking_dataset.json", "w") as f:
    json.dump(pairs, f, indent=2)
print(f"Generated {len(pairs)} pairs")
\end{minted}

% ============================================================
\section{Project 3: Image generation pipeline}

\textbf{Goal:} Build an image generation application with style control and variations.

\textbf{Chapters used:} 7 (diffusion models).

\subsection{Specification}

\begin{enumerate}
  \item Create a pipeline that generates images from text prompts.
  \item Implement style presets (photorealistic, watercolor, digital art, sketch).
  \item Add image-to-image for variations of existing images.
  \item Use ControlNet for edge-guided generation.
  \item Generate a gallery of 20 images across 5 different styles.
\end{enumerate}

\begin{minted}{python}
from diffusers import StableDiffusionPipeline, StableDiffusionImg2ImgPipeline
import torch

class ImageGenerator:
    STYLE_PROMPTS = {
        "photorealistic": "photorealistic, 8k, highly detailed, sharp focus",
        "watercolor": "watercolor painting, soft colors, artistic",
        "digital_art": "digital art, vibrant colors, detailed illustration",
        "sketch": "pencil sketch, black and white, detailed drawing",
        "oil_painting": "oil painting, rich textures, masterpiece",
    }

    def __init__(self, model_id="stabilityai/stable-diffusion-2-1-base"):
        self.pipe = StableDiffusionPipeline.from_pretrained(
            model_id, torch_dtype=torch.float16
        ).to("cuda")

    def generate(self, prompt, style="photorealistic", seed=42, **kwargs):
        styled_prompt = f"{prompt}, {self.STYLE_PROMPTS[style]}"
        generator = torch.Generator("cuda").manual_seed(seed)
        image = self.pipe(
            styled_prompt,
            negative_prompt="blurry, low quality, deformed",
            generator=generator,
            num_inference_steps=30,
            guidance_scale=7.5,
            **kwargs,
        ).images[0]
        return image

gen = ImageGenerator()
for style in ["photorealistic", "watercolor", "digital_art", "sketch"]:
    img = gen.generate("A bustling market in Cotonou", style=style)
    img.save(f"market_{style}.png")
    print(f"Generated: market_{style}.png")
\end{minted}

% ============================================================
\section{Project 4: Multi-agent research assistant}

\textbf{Goal:} Build a multi-agent system where agents collaborate to research, draft, and review a report.

\textbf{Chapters used:} 4 (prompting), 6 (RAG), 9 (agents).

\subsection{Specification}

\begin{enumerate}
  \item \textbf{Researcher agent}: searches for information and compiles notes.
  \item \textbf{Writer agent}: drafts a structured report from the notes.
  \item \textbf{Reviewer agent}: checks for accuracy, coherence, and completeness.
  \item Use LangGraph to orchestrate the workflow with conditional loops.
  \item The reviewer can send the draft back to the writer for revision (max 2 rounds).
\end{enumerate}

\begin{minted}{python}
from langgraph.graph import StateGraph, END
from langchain_google_genai import ChatGoogleGenerativeAI
from typing import TypedDict

class ResearchState(TypedDict):
    topic: str
    research_notes: str
    draft: str
    review: str
    revision_count: int
    final_report: str

llm = ChatGoogleGenerativeAI(model="gemini-2.0-flash", temperature=0.3)

def researcher(state):
    result = llm.invoke(
        f"You are a researcher. Compile detailed notes on: {state['topic']}. "
        f"Include key facts, statistics, and recent developments."
    )
    return {"research_notes": result.content}

def writer(state):
    result = llm.invoke(
        f"You are a technical writer. Write a structured report based on:\n"
        f"{state['research_notes']}\n"
        f"{'Previous review: ' + state.get('review', '') if state.get('review') else ''}"
    )
    return {"draft": result.content,
            "revision_count": state.get("revision_count", 0) + 1}

def reviewer(state):
    result = llm.invoke(
        f"You are a critical reviewer. Review this report for accuracy, "
        f"completeness, and clarity:\n{state['draft']}\n"
        f"Respond with APPROVED if the report is good, or list specific improvements."
    )
    return {"review": result.content}

def should_revise(state):
    if state.get("revision_count", 0) >= 3:
        return "finalize"
    if "APPROVED" in state.get("review", ""):
        return "finalize"
    return "revise"

def finalize(state):
    return {"final_report": state["draft"]}

# Build graph
workflow = StateGraph(ResearchState)
workflow.add_node("research", researcher)
workflow.add_node("write", writer)
workflow.add_node("review", reviewer)
workflow.add_node("finalize", finalize)

workflow.set_entry_point("research")
workflow.add_edge("research", "write")
workflow.add_edge("write", "review")
workflow.add_conditional_edges("review", should_revise,
                                {"revise": "write", "finalize": "finalize"})
workflow.add_edge("finalize", END)

app = workflow.compile()
result = app.invoke({"topic": "The impact of AI on education in Africa",
                      "revision_count": 0})
print(result["final_report"])
\end{minted}

% ============================================================
\section{Project 5: Evaluation benchmark}

\textbf{Goal:} Build a comprehensive evaluation benchmark for comparing LLMs.

\textbf{Chapters used:} 3 (generation), 4 (prompting), 8 (evaluation).

\subsection{Specification}

\begin{enumerate}
  \item Design 50 test cases across 5 categories: factual QA, reasoning, code generation, summarization, and safety.
  \item Evaluate 3 models (GPT-2, TinyLlama, Gemini Flash) on all test cases.
  \item Compute automated metrics: BLEU, ROUGE, perplexity.
  \item Conduct human evaluation on a 20-case subset.
  \item Produce a comparison report with tables and visualizations.
\end{enumerate}

\begin{minted}{python}
import json
from rouge_score import rouge_scorer
import numpy as np

class LLMBenchmark:
    def __init__(self):
        self.scorer = rouge_scorer.RougeScorer(
            ["rouge1", "rouge2", "rougeL"], use_stemmer=True
        )
        self.results = []

    def add_test_case(self, category, question, reference_answer):
        self.results.append({
            "category": category,
            "question": question,
            "reference": reference_answer,
            "model_outputs": {},
        })

    def evaluate_model(self, model_name, generate_fn):
        for case in self.results:
            output = generate_fn(case["question"])
            scores = self.scorer.score(case["reference"], output)
            case["model_outputs"][model_name] = {
                "output": output,
                "rouge1_f1": scores["rouge1"].fmeasure,
                "rouge2_f1": scores["rouge2"].fmeasure,
                "rougeL_f1": scores["rougeL"].fmeasure,
            }

    def summary(self):
        for model in list(self.results[0]["model_outputs"].keys()):
            r1 = np.mean([c["model_outputs"][model]["rouge1_f1"] for c in self.results])
            r2 = np.mean([c["model_outputs"][model]["rouge2_f1"] for c in self.results])
            rl = np.mean([c["model_outputs"][model]["rougeL_f1"] for c in self.results])
            print(f"{model}: ROUGE-1={r1:.3f}, ROUGE-2={r2:.3f}, ROUGE-L={rl:.3f}")

# Usage
bench = LLMBenchmark()
bench.add_test_case("factual", "What is the capital of Benin?", "The capital of Benin is Porto-Novo.")
bench.add_test_case("reasoning", "If A > B and B > C, is A > C?", "Yes, by transitivity, A > C.")
# ... add more test cases
\end{minted}

% ============================================================
\section{Project grading rubric}

\begin{longtable}{p{4cm}p{2cm}p{6.5cm}}
\toprule
\textbf{Criterion} & \textbf{Points} & \textbf{Description} \\
\midrule
Working code & 30 & Code runs without errors on Colab \\
Technical depth & 25 & Correct use of techniques from the course \\
Evaluation & 20 & Quantitative and/or qualitative evaluation \\
Documentation & 15 & Clear notebook with explanations \\
Reflection & 10 & Discussion of limitations and improvements \\
\midrule
\textbf{Total} & \textbf{100} & \\
\bottomrule
\end{longtable}

\begin{exercise}
\begin{enumerate}
  \item Choose one of the five projects and implement it completely.
  \item Write a 500-word project report covering: objective, approach, results, limitations, and future work.
  \item Present your project in a 10-minute demo showing the working application.
  \item Peer-review one other student's project using the grading rubric above.
  \item Propose a sixth capstone project idea that combines at least 3 chapters from this course. Write the specification.
\end{enumerate}
\end{exercise}

\begin{ethicsbox}
Your capstone project will likely be the first generative AI system you build. Before sharing it:
\begin{itemize}
  \item Test for biased or harmful outputs.
  \item Add appropriate disclaimers (``AI-generated content'').
  \item Do not fine-tune on copyrighted or private data without permission.
  \item Consider the social impact: who benefits and who might be harmed?
\end{itemize}
Every AI practitioner is responsible for the systems they build.
\end{ethicsbox}

% ============================================================
\section{Chapter summary}

\begin{itemize}
  \item Five capstone projects cover RAG, fine-tuning, image generation, multi-agent systems, and evaluation.
  \item Each project integrates concepts from multiple course chapters.
  \item Projects are evaluated on working code, technical depth, evaluation, documentation, and reflection.
  \item Build responsibly: test for harm, document limitations, and add appropriate disclaimers.
\end{itemize}
